We sweep $w_E$ (iLQR-Energy) and the scale $q$ of the quadratic state cost (iLQR-Quad) over $\{0.01,\dots,100\}$ with seeds 0--2 (Table~\ref{tab:sweep}, Figs.~\ref{fig:swsr} and~\ref{fig:sweff}); the other weights stay at their defaults.

\begin{table*}[t]\centering\caption{Weight sensitivity: mean over seeds 0--2 at each value of the swept weight ($w_E$ for the energy cost, $q$ for the quadratic cost).}\label{tab:sweep}
\begin{tabular}{lllrrrrr}\toprule
Cost & Task & Metric & 0.01 & 0.1 & 1 & 10 & 100 \\\midrule
Energy & pend. & success & \rhval{sweep_we/ilqr-energy-ours@we=0.01/pendulum/success_rate/mean:2} & \rhval{sweep_we/ilqr-energy-ours@we=0.1/pendulum/success_rate/mean:2} & \rhval{sweep_we/ilqr-energy-ours@we=1/pendulum/success_rate/mean:2} & \rhval{sweep_we/ilqr-energy-ours@we=10/pendulum/success_rate/mean:2} & \rhval{sweep_we/ilqr-energy-ours@we=100/pendulum/success_rate/mean:2} \\
Energy & pend. & effort & \rhval{sweep_we/ilqr-energy-ours@we=0.01/pendulum/control_effort/mean:1} & \rhval{sweep_we/ilqr-energy-ours@we=0.1/pendulum/control_effort/mean:1} & \rhval{sweep_we/ilqr-energy-ours@we=1/pendulum/control_effort/mean:1} & \rhval{sweep_we/ilqr-energy-ours@we=10/pendulum/control_effort/mean:1} & \rhval{sweep_we/ilqr-energy-ours@we=100/pendulum/control_effort/mean:1} \\
Energy & pend. & time (s) & \rhval{sweep_we/ilqr-energy-ours@we=0.01/pendulum/time_to_upright/mean:2} & \rhval{sweep_we/ilqr-energy-ours@we=0.1/pendulum/time_to_upright/mean:2} & \rhval{sweep_we/ilqr-energy-ours@we=1/pendulum/time_to_upright/mean:2} & \rhval{sweep_we/ilqr-energy-ours@we=10/pendulum/time_to_upright/mean:2} & \rhval{sweep_we/ilqr-energy-ours@we=100/pendulum/time_to_upright/mean:2} \\
Energy & cart-p. & success & \rhval{sweep_we/ilqr-energy-ours@we=0.01/cartpole/success_rate/mean:2} & \rhval{sweep_we/ilqr-energy-ours@we=0.1/cartpole/success_rate/mean:2} & \rhval{sweep_we/ilqr-energy-ours@we=1/cartpole/success_rate/mean:2} & \rhval{sweep_we/ilqr-energy-ours@we=10/cartpole/success_rate/mean:2} & \rhval{sweep_we/ilqr-energy-ours@we=100/cartpole/success_rate/mean:2} \\
Energy & cart-p. & effort & \rhval{sweep_we/ilqr-energy-ours@we=0.01/cartpole/control_effort/mean:1} & \rhval{sweep_we/ilqr-energy-ours@we=0.1/cartpole/control_effort/mean:1} & \rhval{sweep_we/ilqr-energy-ours@we=1/cartpole/control_effort/mean:1} & \rhval{sweep_we/ilqr-energy-ours@we=10/cartpole/control_effort/mean:1} & \rhval{sweep_we/ilqr-energy-ours@we=100/cartpole/control_effort/mean:1} \\
Energy & cart-p. & time (s) & \rhval{sweep_we/ilqr-energy-ours@we=0.01/cartpole/time_to_upright/mean:2} & \rhval{sweep_we/ilqr-energy-ours@we=0.1/cartpole/time_to_upright/mean:2} & \rhval{sweep_we/ilqr-energy-ours@we=1/cartpole/time_to_upright/mean:2} & \rhval{sweep_we/ilqr-energy-ours@we=10/cartpole/time_to_upright/mean:2} & \rhval{sweep_we/ilqr-energy-ours@we=100/cartpole/time_to_upright/mean:2} \\
\midrule
Quad & pend. & success & \rhval{sweep_qscale/ilqr-quad@qscale=0.01/pendulum/success_rate/mean:2} & \rhval{sweep_qscale/ilqr-quad@qscale=0.1/pendulum/success_rate/mean:2} & \rhval{sweep_qscale/ilqr-quad@qscale=1/pendulum/success_rate/mean:2} & \rhval{sweep_qscale/ilqr-quad@qscale=10/pendulum/success_rate/mean:2} & \rhval{sweep_qscale/ilqr-quad@qscale=100/pendulum/success_rate/mean:2} \\
Quad & pend. & effort & \rhval{sweep_qscale/ilqr-quad@qscale=0.01/pendulum/control_effort/mean:1} & \rhval{sweep_qscale/ilqr-quad@qscale=0.1/pendulum/control_effort/mean:1} & \rhval{sweep_qscale/ilqr-quad@qscale=1/pendulum/control_effort/mean:1} & \rhval{sweep_qscale/ilqr-quad@qscale=10/pendulum/control_effort/mean:1} & \rhval{sweep_qscale/ilqr-quad@qscale=100/pendulum/control_effort/mean:1} \\
Quad & pend. & time (s) & \rhval{sweep_qscale/ilqr-quad@qscale=0.01/pendulum/time_to_upright/mean:2} & \rhval{sweep_qscale/ilqr-quad@qscale=0.1/pendulum/time_to_upright/mean:2} & \rhval{sweep_qscale/ilqr-quad@qscale=1/pendulum/time_to_upright/mean:2} & \rhval{sweep_qscale/ilqr-quad@qscale=10/pendulum/time_to_upright/mean:2} & \rhval{sweep_qscale/ilqr-quad@qscale=100/pendulum/time_to_upright/mean:2} \\
Quad & cart-p. & success & \rhval{sweep_qscale/ilqr-quad@qscale=0.01/cartpole/success_rate/mean:2} & \rhval{sweep_qscale/ilqr-quad@qscale=0.1/cartpole/success_rate/mean:2} & \rhval{sweep_qscale/ilqr-quad@qscale=1/cartpole/success_rate/mean:2} & \rhval{sweep_qscale/ilqr-quad@qscale=10/cartpole/success_rate/mean:2} & \rhval{sweep_qscale/ilqr-quad@qscale=100/cartpole/success_rate/mean:2} \\
Quad & cart-p. & effort & \rhval{sweep_qscale/ilqr-quad@qscale=0.01/cartpole/control_effort/mean:1} & \rhval{sweep_qscale/ilqr-quad@qscale=0.1/cartpole/control_effort/mean:1} & \rhval{sweep_qscale/ilqr-quad@qscale=1/cartpole/control_effort/mean:1} & \rhval{sweep_qscale/ilqr-quad@qscale=10/cartpole/control_effort/mean:1} & \rhval{sweep_qscale/ilqr-quad@qscale=100/cartpole/control_effort/mean:1} \\
Quad & cart-p. & time (s) & \rhval{sweep_qscale/ilqr-quad@qscale=0.01/cartpole/time_to_upright/mean:2} & \rhval{sweep_qscale/ilqr-quad@qscale=0.1/cartpole/time_to_upright/mean:2} & \rhval{sweep_qscale/ilqr-quad@qscale=1/cartpole/time_to_upright/mean:2} & \rhval{sweep_qscale/ilqr-quad@qscale=10/cartpole/time_to_upright/mean:2} & \rhval{sweep_qscale/ilqr-quad@qscale=100/cartpole/time_to_upright/mean:2} \\
\bottomrule
\end{tabular}
\end{table*}

\begin{figure}[t]\centering\includegraphics[width=\columnwidth]{sweep_success.pdf}\caption{Success rate vs.\ the swept weight ($w_E$ for iLQR-Energy, $q$ for iLQR-Quad), per task. Mean $\pm$ one sample std over 3 seeds, log $x$-axis.}\label{fig:swsr}\end{figure}
\begin{figure}[t]\centering\includegraphics[width=\columnwidth]{sweep_effort.pdf}\caption{Control effort vs.\ the swept weight, per task. Mean $\pm$ one sample std over 3 seeds, log $x$-axis.}\label{fig:sweff}\end{figure}

\textbf{H3 (less sensitive to weight): refuted.} The registered test is the spread of the success rate across weights, and it is the same for the two costs on each task: across the five weights the mean success runs from \rhval{sweep_we/ilqr-energy-ours@we=0.01/pendulum/success_rate/mean:2} to \rhval{sweep_we/ilqr-energy-ours@we=100/pendulum/success_rate/mean:2} for the energy cost and from \rhval{sweep_qscale/ilqr-quad@qscale=0.01/pendulum/success_rate/mean:2} to \rhval{sweep_qscale/ilqr-quad@qscale=100/pendulum/success_rate/mean:2} for the quadratic cost on the pendulum, and from \rhval{sweep_we/ilqr-energy-ours@we=0.1/cartpole/success_rate/mean:2} to \rhval{sweep_we/ilqr-energy-ours@we=100/cartpole/success_rate/mean:2} and from \rhval{sweep_qscale/ilqr-quad@qscale=0.01/cartpole/success_rate/mean:2} to \rhval{sweep_qscale/ilqr-quad@qscale=100/cartpole/success_rate/mean:2} on the cart-pole. Effort tells a mixed story: the energy cost's effort varies more than the quadratic cost's on the pendulum (lowest to highest mean \rhval{sweep_we/ilqr-energy-ours@we=100/pendulum/control_effort/mean:1} to \rhval{sweep_we/ilqr-energy-ours@we=0.01/pendulum/control_effort/mean:1} vs \rhval{sweep_qscale/ilqr-quad@qscale=100/pendulum/control_effort/mean:1} to \rhval{sweep_qscale/ilqr-quad@qscale=1/pendulum/control_effort/mean:1}) and less on the cart-pole (\rhval{sweep_we/ilqr-energy-ours@we=1/cartpole/control_effort/mean:1} to \rhval{sweep_we/ilqr-energy-ours@we=100/cartpole/control_effort/mean:1} vs \rhval{sweep_qscale/ilqr-quad@qscale=0.01/cartpole/control_effort/mean:1} to \rhval{sweep_qscale/ilqr-quad@qscale=100/cartpole/control_effort/mean:1}). The two weights are not commensurable ($w_E$ multiplies a squared energy, $q$ a squared state error), so a ``same decade'' comparison is only indicative.

\textbf{Success depends on the weight.} Both costs improve with larger weights on the pendulum: $w_E\ge 10$ and $q=100$ give success \rhval{sweep_we/ilqr-energy-ours@we=10/pendulum/success_rate/mean:2} in the 3 seeds, above the default-weight values. With the best swept weight iLQR-Quad is also at \rhval{sweep_qscale/ilqr-quad@qscale=100/pendulum/success_rate/mean:2}, which removes the H1 gap.

\textbf{Effort across the sweep.} Table~\ref{tab:paired} compares iLQR-Quad at every $q$ with iLQR-Energy at its default $w_E=1$, paired over seeds 0--2. On the pendulum the quadratic cost uses more effort at every $q$ (by \rhval{cmp/pair_q100/ilqr-energy-ours/pendulum/control_effort/delta:1} to \rhval{cmp/pair_q1/ilqr-energy-ours/pendulum/control_effort/delta:1}). On the cart-pole the difference is \rhval{cmp/pair_q1/ilqr-energy-ours/cartpole/control_effort/delta:1} at the default $q=1$ and \rhval{cmp/pair_q10/ilqr-energy-ours/cartpole/control_effort/delta:1} to \rhval{cmp/pair_q100/ilqr-energy-ours/cartpole/control_effort/delta:1} for larger $q$, but only \rhval{cmp/pair_q0.01/ilqr-energy-ours/cartpole/control_effort/delta:1} and \rhval{cmp/pair_q0.1/ilqr-energy-ours/cartpole/control_effort/delta:1} for $q=0.01$ and $q=0.1$, with paired $p$ of \rhval{cmp/pair_q0.01/ilqr-energy-ours/cartpole/control_effort/paired_p:3} and \rhval{cmp/pair_q0.1/ilqr-energy-ours/cartpole/control_effort/paired_p:3} at $n=3$. The reverse reading matters too: iLQR-Energy at $w_E=0.01$, 0.1 and 100 (\rhval{sweep_we/ilqr-energy-ours@we=0.01/cartpole/control_effort/mean:1}, \rhval{sweep_we/ilqr-energy-ours@we=0.1/cartpole/control_effort/mean:1}, \rhval{sweep_we/ilqr-energy-ours@we=100/cartpole/control_effort/mean:1}) is not below iLQR-Quad at $q\le 0.1$ (\rhval{sweep_qscale/ilqr-quad@qscale=0.01/cartpole/control_effort/mean:1}, \rhval{sweep_qscale/ilqr-quad@qscale=0.1/cartpole/control_effort/mean:1}). So on the cart-pole the energy cost is nominally cheaper than the re-weighted quadratic cost only at $w_E=1$ and $w_E=10$ (\rhval{sweep_we/ilqr-energy-ours@we=1/cartpole/control_effort/mean:1} and \rhval{sweep_we/ilqr-energy-ours@we=10/cartpole/control_effort/mean:1}), by a margin that 3 seeds do not clearly resolve.

\begin{table}[t]\centering\caption{Effort of iLQR-Quad at each $q$ minus effort of iLQR-Energy at $w_E=1$ ($\Delta$ effort, positive: energy cost is cheaper), paired $t$-test over seeds 0--2, and the success rate of iLQR-Quad at that $q$.}\label{tab:paired}
\setlength{\tabcolsep}{5pt}\begin{tabular}{llrrrrr}\toprule
Task & & $q{=}0.01$ & 0.1 & 1 & 10 & 100 \\\midrule
pend. & $\Delta$ effort & 4.2 & 4.2 & 4.7 & 4.4 & 2.8 \\
 & paired $p$ & 0.059 & 0.057 & 0.047 & $<$0.001 & 0.044 \\
 & Quad success & 0.83 & 0.83 & 0.83 & 0.92 & 1.00 \\
cart-p. & $\Delta$ effort & 6.7 & 6.9 & 31.5 & 55.4 & 59.1 \\
 & paired $p$ & 0.046 & 0.090 & 0.027 & 0.003 & 0.008 \\
 & Quad success & 0.98 & 1.00 & 1.00 & 1.00 & 1.00 \\
\bottomrule
\end{tabular}
\end{table}

\textbf{Tuned comparison.} To compare the costs at equal tuning budget we pick, for each cost and task, the swept weight with the highest mean success on seeds 0--2 (ties broken by lowest mean effort) and add seeds 3--4 at that weight (Table~\ref{tab:tuned}). This picks $w_E=100$ and $q=100$ on the pendulum, and $w_E=1$ and $q=0.1$ on the cart-pole. The rule was fixed after seeing the sweep, so this is a post-hoc analysis. Success is \rhval{tuned_all/ilqr-energy-ours/pendulum/success_rate/mean:2} vs \rhval{tuned_all/ilqr-quad/pendulum/success_rate/mean:2} on the pendulum and \rhval{tuned_all/ilqr-energy-ours/cartpole/success_rate/mean:2} vs \rhval{tuned_all/ilqr-quad/cartpole/success_rate/mean:2} on the cart-pole: no success advantage remains. Effort stays lower for the energy cost, \rhval{tuned_all/ilqr-energy-ours/pendulum/control_effort/mean:1} vs \rhval{tuned_all/ilqr-quad/pendulum/control_effort/mean:1} on the pendulum and \rhval{tuned_all/ilqr-energy-ours/cartpole/control_effort/mean:1} vs \rhval{tuned_all/ilqr-quad/cartpole/control_effort/mean:1} on the cart-pole (paired $p$ of \rhval{cmp/tuned_all/ilqr-energy-ours/pendulum/control_effort/paired_p:4} and \rhval{cmp/tuned_all/ilqr-energy-ours/cartpole/control_effort/paired_p:3} at $n=5$). The 5 seeds include the 3 used for selection, for both costs alike; on the two held-out seeds alone the means are \rhval{tuned/ilqr-energy-ours/pendulum/control_effort/mean:1} vs \rhval{tuned/ilqr-quad/pendulum/control_effort/mean:1} and \rhval{tuned/ilqr-energy-ours/cartpole/control_effort/mean:1} vs \rhval{tuned/ilqr-quad/cartpole/control_effort/mean:1}, with no test at $n=2$. Time to upright is not distinguishable (\rhval{tuned_all/ilqr-energy-ours/pendulum/time_to_upright/mean:2} vs \rhval{tuned_all/ilqr-quad/pendulum/time_to_upright/mean:2}\,s; \rhval{tuned_all/ilqr-energy-ours/cartpole/time_to_upright/mean:2} vs \rhval{tuned_all/ilqr-quad/cartpole/time_to_upright/mean:2}\,s). Only one weight per cost was tuned; $r$, $Q_f$ and the horizon were not.

\begin{table}[t]\centering\caption{Tuned comparison: each cost at the weight picked on seeds 0--2 (highest success, ties by lowest effort). Mean $\pm$ sample std over seeds 0--4; last column: mean effort on the held-out seeds 3--4 only. Paired $t$-test between the two costs over the 5 seeds (--: no variance).}\label{tab:tuned}
\setlength{\tabcolsep}{3.5pt}\begin{tabular}{llrrrr}\toprule
Task & Cost (weight) & success & effort & time (s) & eff.\ 3--4 \\\midrule
pend. & Energy ($w_E$=100) & 0.99$\pm$0.02 & 10.9$\pm$1.5 & 2.29$\pm$0.18 & 12.0 \\
pend. & Quad ($q$=100) & 1.00$\pm$0.00 & 14.4$\pm$1.2 & 2.26$\pm$0.15 & 15.3 \\
 & paired $p$ ($n{=}5$) & 0.374 & $<$0.001 & 0.374 & \\
\midrule
cart-p. & Energy ($w_E$=1) & 1.00$\pm$0.00 & 42.0$\pm$6.0 & 1.79$\pm$0.06 & 40.1 \\
cart-p. & Quad ($q$=0.1) & 1.00$\pm$0.00 & 48.6$\pm$5.5 & 1.75$\pm$0.08 & 46.2 \\
 & paired $p$ ($n{=}5$) & -- & 0.009 & 0.103 & \\
\bottomrule
\end{tabular}
\end{table}
